\begin{abstract}
We formulate a generative model for directional data using spin coherent states, isotropic open-system diffusion, and trainable quantum channels. Classical distributions on the sphere are encoded as mixed spin-$j$ states, whose irreducible multipoles evolve with the spherical heat spectrum $\exp[-D\ell(\ell+1)t]$ up to rank $2j$. We derive the encoding--diffusion correspondence and explicitly distinguish the input density from its finite-resolution Husimi readout. Exact density-matrix simulations then separate representational resolution, initialization, supervision, and optimization budget. In a 40-run paired study at $j=5/2$, an unchanged 162-parameter reverse architecture recovers both encoded target modes in five of five seeds with final-only fidelity training at 100 epochs, compared with one of five under path supervision. A hybrid fidelity--multipole objective selected on separate validation seeds reduces mean high-rank squared error by \HybridRankReduction\% and mean trace distance by \HybridTraceReduction\% on five new confirmation seeds; both objectives recover both modes in all five runs. A two-spin pilot demonstrates the product-phase-space construction but performs worse than correlation-aware classical baselines. These results establish a reproducible finite-dimensional prototype and expose training and representation limitations. They do not establish quantum advantage, hardware efficiency, or a benefit of diffusion supervision over direct state preparation.
\end{abstract}
